% Zakończenie pracy — liczby zgodne z reports/ (stan 2026-07-07).

\section*{Zakończenie}
\addcontentsline{toc}{section}{Zakończenie}

Praca postawiła pytanie, czy ocenę ryzyka kredytowego warto traktować jako
proces rozwijający się w~czasie, a~nie jednorazowy werdykt --- i~odpowiedziała
na nie konstrukcyjnie oraz empirycznie. Zbudowano kompletny system
eksperymentalny (React / .NET~8 / Flask / PostgreSQL), w~którym pięć
skalibrowanych modeli ocenia tę samą ekspozycję na czterech przesuwanych oknach
historii płatniczej, a~trajektoria PD zasila regułę wczesnego ostrzegania
z~progami optymalnymi kosztowo i~wyjaśnieniami predykcji.

Wyniki empiryczne układają się w~spójny obraz. Skrócenie okna obserwacji do
trzech miesięcy kosztuje mniej niż 1~pp AUC (H1). Reguła monitorująca nie
przewyższa statycznej pod względem czułości przy stałym budżecie fałszywych
alarmów --- z~wyjątkiem modelu sekwencyjnego --- ale oferuje około dwóch okien
wyprzedzenia i~kilkadziesiąt wykryć na model niedostępnych ocenie jednorazowej
(H2, potwierdzona częściowo). Wszystkie modele przechodzą audyt sprawiedliwości
z~wielokrotnym zapasem, a~eksperyment ablacyjny pokazał, że usunięcie atrybutu
chronionego z~cech nie kosztuje nic predykcyjnie i~poprawia parytet (H3).
Najlepszym pojedynczym modelem jest CatBoost (AUC 0,779); przewaga metod
drzewiastych nad LSTM na danych tabelarycznych potwierdza literaturę, lecz to
LSTM --- jedyny model bez cech demograficznych i~jedyny architektonicznie
sekwencyjny --- okazuje się najbliższy parytetu i~jako jedyny zyskuje na regule
monitorującej także w~czułości.

Kierunki dalszych badań wynikają wprost z~ograniczeń: walidacja schematu na
prawdziwym panelu podłużnym z~etykietami przesuwanymi w~czasie; złożenie
predykcji w~zespół typu stacking z~protokołem out-of-fold; rozszerzenie audytu
sprawiedliwości na pozostałe atrybuty i~analiza mitygacji progowej; wreszcie
kalibracja modelu kosztów FN/FP na rzeczywistych stratach portfelowych, od
której zależy bilans reguły monitorującej. Niezależnie od tych rozszerzeń,
zasadniczy wniosek pracy pozostaje: monitoring trajektorii PD jest praktycznie
wykonalnym, audytowalnym i~regulacyjnie umocowanym uzupełnieniem oceny
statycznej --- wartym wdrożenia tam, gdzie koszt spóźnionej reakcji przewyższa
koszt dodatkowego przeglądu.
